% reference_exam.tex
% Structural reference for the lecture-materials-assistant exam artifact.
% All course-specific content replaced with bracketed placeholders.
%
% PORTED 2026-05-13 to the new Sp26 production format:
%   - Charter + Helvetica + Inconsolata font stack
%   - \metastrip{course}{term}{exam}{points} metadata strip
%   - \examsection{NAME} small-caps section heading
%   - \correctchoice / \wrongchoice semantic MC markup (color in answers mode)
%   - \doublerule letterpress double rule under the title
%   - \buildtime pdfTeX-primitive timestamp in left foot
%   - \writelines{N} ruled response lines for short-answer questions
%   - Even-page padding for duplex printing
%   - Editorial \setlist metrics
% Preserves per-student layer (\studentname / \studentserial / \issuedto) and
% the answers-mode bridge for pa-exam-build compatibility.
%
% FILE NAMING: [course_num]-exam-[n]-[term]
%   term format : sp/fa/su + 2-digit year  (e.g. sp26, fa25)
%   parallel sections : append -a, -b      (e.g. 326-exam-1-sp26-a.tex)
%
% STUDENT COPY : default (\answersfalse)            → student PDF
% ANSWER KEY   : \def\answersmode{1} on cmdline     → key PDF (or \answerstrue inline)

\documentclass[11pt]{article}
\usepackage[margin=0.9in,top=0.7in,bottom=0.85in,headheight=14pt,headsep=14pt]{geometry}
\usepackage[T1]{fontenc}
\usepackage[utf8]{inputenc}

% --- Print-optimized font stack ---
\usepackage{charter}                          % Bitstream Charter body
\usepackage[scaled=0.92]{helvet}              % Helvetica-clone sans (metadata)
\usepackage[scaled=0.95,varl]{inconsolata}    % Inconsolata mono, varl alt
\renewcommand{\familydefault}{\rmdefault}
\linespread{1.06}
\frenchspacing

\usepackage{enumitem}
\usepackage{needspace}
\usepackage[hyphens]{url}
\usepackage[hidelinks]{hyperref}
\usepackage{listings}
\usepackage{microtype}
\usepackage{fancyhdr}
\usepackage{pifont}                           % Zapf Dingbats — manicule (\ding{43})
\usepackage{xcolor}
\definecolor{keyred}{HTML}{B71C1C}
\definecolor{keygreen}{HTML}{1B5E20}
\definecolor{keygrey}{HTML}{616161}

% --- MC choice color: green-bold for correct, grey for distractors.
% Gated by \ifanswers so student PDFs render identically to the pre-color source.
\newcommand{\correctchoice}[1]{\ifanswers\textcolor{keygreen}{\textbf{#1}}\else #1\fi}
\newcommand{\wrongchoice}[1]{\ifanswers\textcolor{keygrey}{#1}\else #1\fi}

% --- Build-time stamp (pdfTeX primitives) ---
% \time is minutes since midnight. Split into HH:MM by integer division.
\newcount\stamphour
\newcount\stampminute
\stamphour=\time
\divide\stamphour by 60
\stampminute=\time
\count255=\stamphour
\multiply\count255 by 60
\advance\stampminute by -\count255
\newcommand{\padtwo}[1]{\ifnum#1<10 0\fi\the#1}
\newcommand{\buildtime}{\the\year-\padtwo{\month}-\padtwo{\day}~\padtwo{\stamphour}:\padtwo{\stampminute}}

% --- Answer key toggle ---
% pa-exam-build injects \def\answersmode{1} on the cmdline for the key build;
% bridge it to \ifanswers. LaTeX-native \answerstrue still works for direct edits.
\newif\ifanswers
\answersfalse
\makeatletter
\@ifundefined{answersmode}{}{\answerstrue}
\makeatother

% --- Exam serial (variant-level; per-file fingerprint set at build time by pa-exam-build) ---
\providecommand{\examserial}{XXXXXXXX}

% Per-student exam ID macros (see plans/specs/2026-05-13-per-student-exam-id-design).
% Defaults to empty for backward-compat with variant-only builds.
%
% NB: \@ifundefined+\def is used here (not \providecommand) so the default empty
% value produces a non-\long macro, which \ifx-compares equal to \@empty below.
% \providecommand{}{} produces \long macro:-> which is NOT \ifx-equal to \@empty,
% causing the conditional to always render the "· ID" / "Issued to:" branches.
\makeatletter
\@ifundefined{studentname}{\def\studentname{}}{}
\@ifundefined{studentserial}{\def\studentserial{}}{}
\makeatother

% \issuedto --- small grey "Issued to: <name>" line above the NAME block.
% Renders only if \studentname is set.
\makeatletter
\newcommand{\issuedto}{%
  \ifx\studentname\@empty\else
    \noindent{\sffamily\footnotesize\color{black!55}Issued to: \studentname}%
    \par\vspace{4pt}%
  \fi%
}
\makeatother

% --- Running header + italic foot ---
% Header: italic small charter — feels editorial, not corporate.
% Foot: italic "p. N" — old-school typographic pagination.
% Right foot: serial + optional per-student ID, both monospaced.
\pagestyle{fancy}
\fancyhf{}
\fancyhead[L]{\textit{\small [COURSE CODE] · [TERM]}}
\fancyhead[R]{\textit{\small [Exam Name]}}
\fancyfoot[L]{\textit{\footnotesize Generated \buildtime}}
\fancyfoot[C]{\textit{\small p.~\thepage}}
\makeatletter
\fancyfoot[R]{\textit{\footnotesize Serial \texttt{\examserial}%
  \ifx\studentserial\@empty\else
    \ \textperiodcentered\ ID \texttt{\studentserial}%
  \fi}}
\makeatother
\renewcommand{\headrulewidth}{0.4pt}
\renewcommand{\footrulewidth}{0pt}

% --- Letterpress double rule (thick + thin, 2pt apart) ---
% Used under the title block. Centuries-old letterpress convention:
% the double rule signals "title above".
\newcommand{\doublerule}{%
  \noindent\rule{\linewidth}{1.0pt}\par\vspace{1.5pt}%
  \noindent\rule{\linewidth}{0.3pt}\par}

% --- Drafting-style metadata strip ---
% Two-column key/value table in small caps for labels, charter for values.
% Hairline rule above and below; mimics an engineer's drawing title block
% without the heavy chrome.
\newcommand{\metastrip}[4]{%
  \par\noindent\rule{\linewidth}{0.3pt}\par\vspace{2pt}%
  \noindent\begin{minipage}{\linewidth}%
  \begin{tabbing}%
  \hspace{6em}\=\hspace{16em}\=\hspace{6em}\=\kill%
  \textsc{course}~~\>#1\>\textsc{term}~~\>#2\\[1pt]%
  \textsc{exam}~~\>#3\>\textsc{points}~~\>#4%
  \end{tabbing}%
  \end{minipage}\par\vspace{-4pt}%
  \noindent\rule{\linewidth}{0.3pt}\par}

% --- Section-heading command for question groups ---
% Small caps + letterspaced for editorial gravitas; thin rule above.
\newcommand{\examsection}[1]{%
  \par\vspace{1.4em}%
  \noindent\rule{\linewidth}{0.3pt}\par\vspace{2pt}%
  \noindent{\textsc{\large #1}}%
  \par\vspace{0.4em}}

% --- Ornamental section break (optional, between MC and SA) ---
\newcommand{\examornament}{%
  \par\vspace{0.6em}\begin{center}\large $*$\quad$*$\quad$*$\end{center}\par}

% --- List metrics: editorial proportions, tightened ---
\setlist[enumerate,1]{leftmargin=2.4em, itemsep=0.45em, topsep=0.45em, parsep=0pt, label=\textbf{\arabic*.}}
\setlist[enumerate,2]{leftmargin=1.8em, itemsep=0pt, topsep=0.1em, parsep=0pt, label=(\alph*)}

% --- Code blocks: hairline-rule top + bottom, no fill, no line numbers ---
% Short exam snippets don't need line-number anchors; the top/bottom hair-rules
% read as a typographic "code" treatment without competing with a number column.
\lstset{
  basicstyle=\ttfamily\small,
  breaklines=true,
  breakatwhitespace=true,
  columns=fullflexible,
  showstringspaces=false,
  numbers=none,
  tabsize=3,
  aboveskip=4mm,
  belowskip=4mm,
  frame=tb,
  framerule=0.4pt,
  framesep=6pt,
  xleftmargin=12pt,
  xrightmargin=12pt,
}

% --- Ruled writing lines for short-answer responses ---
\newcounter{wl}
\newcommand{\writelines}[1]{%
  \par\vspace{0.4em}%
  \setcounter{wl}{0}%
  \loop\ifnum\value{wl}<#1%
    \noindent\rule[0.45em]{\linewidth}{0.3pt}\par\vspace{0.30in}%
    \stepcounter{wl}%
  \repeat}

% --- Even-page padding for duplex printing ---
% Two-sided printing reality: an odd-page exam means the back of the last
% leaf is the front of the NEXT student's exam — bad. Pad with a blank
% verso when total page count is odd.
\AtEndDocument{%
  \clearpage%
  \ifodd\value{page}\else%
    \thispagestyle{plain}%
    \null\vfill%
    \begin{center}\textit{\small This page intentionally left blank.}\end{center}%
    \vfill\null%
  \fi%
}

\setlength{\emergencystretch}{2em}
\sloppy

% --- Key banner activation: when in answers mode, flip the right header. ---
\ifanswers\fancyhead[R]{\textit{\small [Exam Name] --- \textcolor{keyred}{\textbf{KEY}}}}\fi

\date{}

\begin{document}

% -----------------------------------------------------------------------
% TITLE BLOCK
% -----------------------------------------------------------------------

\metastrip{[COURSE CODE]}{[TERM]}{[Exam Name]}{[N]}
\par\vspace{0.4em}
\begin{center}{\Large\itshape [Exam Name]}\end{center}
\par\vspace{-0.2em}
\doublerule
\par\vspace{0.9em}

\par\vspace{6pt}
\noindent{\sffamily\bfseries\small\ding{43}~~PRINT CLEARLY. UNNAMED EXAMS CANNOT BE RETURNED OR GRADED.}
\par\vspace{12pt}

\issuedto
\noindent{\sffamily\bfseries\large NAME:}~\rule[-3pt]{13.5cm}{0.8pt}
\par\vspace{14pt}
\noindent{\sffamily\bfseries\large STUDENT ID\#:}~\rule[-3pt]{5.5cm}{0.8pt}\hfill{\sffamily\bfseries\large DATE:}~\rule[-3pt]{4cm}{0.8pt}
\par\vspace{6pt}

\ifanswers\par\vspace{0.4em}\noindent\colorbox{keyred}{\parbox{\dimexpr\textwidth-2\fboxsep\relax}{\centering\color{white}\sffamily\bfseries\large $\star$~~ANSWER KEY~~---~~NOT FOR DISTRIBUTION~~$\star$}}\par\vspace{0.4em}\fi

% -----------------------------------------------------------------------
% DIRECTIONS
% -----------------------------------------------------------------------

\par\vspace{0.6em}
\noindent\ding{43}~\textbf{Directions.}~\textit{Write on the exam paper. No electronic devices allowed. Answer completely but as briefly as possible. All answers must be in your own words.}
\par\vspace{0.4em}

% -----------------------------------------------------------------------
% MULTIPLE CHOICE
% -----------------------------------------------------------------------

\examsection{Multiple Choice}
\begin{enumerate}

% --- mc: standard 4-option (use \correctchoice / \wrongchoice for semantic markup) ---
\item \textit{([N] pts)}~Which of the following best describes [concept]?
  \begin{enumerate}[label=(\alph*)]
    \item \wrongchoice{[plausible distractor]}
    \item \wrongchoice{[plausible distractor]}
    \item \correctchoice{[correct answer]}
    \item \wrongchoice{[plausible distractor]}
  \end{enumerate}
  \ifanswers \textcolor{keyred}{\textbf{Answer:} c} \fi

% --- tf: true/false ---
\item \textit{([N] pts)}~\textsc{T~/~F.}~[Statement about a concept that is either true or false.]
  \ifanswers \textcolor{keyred}{\textbf{Answer:} True} \fi

% --- code: code-interpretation T/F ---
% Show a complete, self-contained implementation (~15 lines max).
% Incorrect implementations must have a subtle, realistic bug. Don't place
% code questions back-to-back.
\item \textit{([N] pts)}~\textsc{T~/~F.}~The following implementation is correct:
\begin{lstlisting}
// Show a complete, self-contained implementation.
// Correct or subtly broken -- make the bug realistic.
void example(void)
{
    // ...
}
\end{lstlisting}
  \ifanswers \textcolor{keyred}{\textbf{Answer:} False --- [describe the specific bug]} \fi

\end{enumerate}

% -----------------------------------------------------------------------
% SHORT ANSWER / ESSAY  (optional — comment out the whole section if MC-only)
% -----------------------------------------------------------------------

\examsection{Short Answer}
\begin{enumerate}[resume]

\item \needspace{2.5in} \textit{([N] pts)}~[Synthesis / compare-contrast / analyze-scenario prompt. Should require thinking — not answerable from a single definition.]
  \ifanswers \par\smallskip \textcolor{keyred}{\textit{Key:}} [Model answer usable as a grading rubric.]
  \par\smallskip \textsc{\footnotesize Coverage:}~{\footnotesize [lecture S#] $\cdot$ [textbook §] $\cdot$ [reading list reference]}
  \par\smallskip \textbf{Gradescope rubric} \textit{(additive, [N] pts):}
  \begin{itemize}[noitemsep, leftmargin=1.5em, topsep=2pt]
    \item \textbf{\textcolor{keygreen}{+[N]}} \quad [criterion satisfied] \textit{([rubric-category])}
    \item \textbf{\textcolor{keygrey}{+0}} \quad [criterion not met] \textit{([rubric-category])}
    \item \textbf{\textcolor{keygrey}{+0}} \quad Blank / off-topic \textit{(catchall)}
  \end{itemize}
  \else \writelines{16} \fi

\item \needspace{2.5in} \textit{([N] pts)}~[Second synthesis or code-interpretation question.]
  \ifanswers \par\smallskip \textcolor{keyred}{\textit{Key:}} [Model answer.]
  \par\smallskip \textsc{\footnotesize Coverage:}~{\footnotesize [lecture S#] $\cdot$ [textbook §] $\cdot$ [reading list reference]}
  \par\smallskip \textbf{Gradescope rubric} \textit{(additive, [N] pts):}
  \begin{itemize}[noitemsep, leftmargin=1.5em, topsep=2pt]
    \item \textbf{\textcolor{keygreen}{+[N]}} \quad [criterion] \textit{([rubric-category])}
    \item \textbf{\textcolor{keygrey}{+0}} \quad Blank / off-topic \textit{(catchall)}
  \end{itemize}
  \else \writelines{16} \fi

\end{enumerate}

\par\vfill
\begin{center}\small\textit{End of exam.}\end{center}

\end{document}
